
\documentclass{article}
\usepackage{colortbl}
\usepackage{makecell}
\usepackage{multirow}
\usepackage{supertabular}

\begin{document}

\newcounter{utterance}

\centering \large Interaction Transcript for game `referencegame', experiment `line\_grids\_columns', episode 3 with NEW\_Reflect\_Agent.
\vspace{24pt}

{ \footnotesize  \setcounter{utterance}{1}
\setlength{\tabcolsep}{0pt}
\begin{supertabular}{c@{$\;$}|p{.15\linewidth}@{}p{.15\linewidth}p{.15\linewidth}p{.15\linewidth}p{.15\linewidth}p{.15\linewidth}}
    \# & $\;$A & \multicolumn{4}{c}{Game Master} & $\;\:$B\\
    \hline

    \theutterance \stepcounter{utterance}  
    & & \multicolumn{4}{p{0.6\linewidth}}{
        \cellcolor[rgb]{0.9,0.9,0.9}{
            \makecell[{{p{\linewidth}}}]{
                \texttt{\tiny{[P1$\langle$GM]}}
                \texttt{You are given three grids, where each of them is 5 by 5 in size.} \\
\texttt{Grids have empty cells marked with "▢" and filled cells marked with "X".} \\
\texttt{Your task is to generate a referring expression that best describes the target grid while distinguishing it from the two other distractor grids.} \\
\texttt{The first grid is the target grid, and the following two grids are the distractors.} \\
\\ 
\texttt{Target grid:} \\
\\ 
\texttt{X X X □ □} \\
\texttt{X X X □ □} \\
\texttt{X X X □ □} \\
\texttt{X X X □ □} \\
\texttt{X X X □ □} \\
\\ 
\texttt{Distractor grid 1:} \\
\\ 
\texttt{□ X X X □} \\
\texttt{□ X X X □} \\
\texttt{□ X X X □} \\
\texttt{□ X X X □} \\
\texttt{□ X X X □} \\
\\ 
\texttt{Distractor grid 2:} \\
\\ 
\texttt{□ X X □ X} \\
\texttt{□ X X □ X} \\
\texttt{□ X X □ X} \\
\texttt{□ X X □ X} \\
\texttt{□ X X □ X} \\
\\ 
\texttt{Instruction: Describe the target grid.} \\
\texttt{Generate the referring expression starting with the tag "Expression: " for the given target grid. Omit any other text.} \\
            }
        }
    }
    & \\ \\

    \theutterance \stepcounter{utterance}  
    & \multicolumn{4}{p{0.45\linewidth}}{
        \cellcolor[rgb]{0.95,0.95,0.95}{
            \makecell[{{p{\linewidth}}}]{
                \texttt{\tiny{[P1$|$internal]}}
                \texttt{Reflections from previous episodes:} \\
\texttt{{-}} \\
\texttt{ Failed. My expression didn't distinguish the target because I focused only on the left column of X's, which is also present in distractor 2. I should have specified the exact pattern: a vertical column of X's on the left, then a vertical column of empty cells, then two vertical columns of X's. Next time, I must describe the full unique column pattern, not just one part.} \\
\texttt{{-}} \\
\texttt{ I succeeded. The target grid has a solid vertical block of X's on the right three columns, with the left two columns empty. The distractors have different patterns: one has a vertical empty strip in the middle, the other has empty cells at the leftmost and rightmost columns. My expression likely focused on the rightmost three columns being completely filled, which uniquely identified the target. Next time, I would similarly look for the largest contiguous filled region and its position relative to the grid edges.} \\
\texttt{{-}} \\
\texttt{ Expression: A vertical stripe of X's in the middle three columns, with empty cells on the leftmost and rightmost columns.} \\
\\ 
\texttt{Reflection: Succeeded. I identified the key distinguishing feature: the target has a solid vertical block of X's occupying columns 2–4, while both distractors have the block shifted left or right, leaving two empty columns on one side. I described it clearly and concisely. Next time, I would use the same approach—focus on the shape and position of the filled region relative to the grid edges.} \\
            }
        }
    }
    & & \\ \\

    \theutterance \stepcounter{utterance}  
    & \multicolumn{4}{p{0.45\linewidth}}{
        \cellcolor[rgb]{0.95,0.95,0.95}{
            \makecell[{{p{\linewidth}}}]{
                \texttt{\tiny{[P1$|$internal]}}
                \texttt{{[}attempt 1 – valid{]}} \\
\texttt{Expression: A vertical block of X's in the left three columns, with empty cells in the right two columns.} \\
            }
        }
    }
    & & \\ \\

    \theutterance \stepcounter{utterance}  
    & \multicolumn{4}{p{0.6\linewidth}}{
        \cellcolor[rgb]{0.9,0.9,0.9}{
            \makecell[{{p{\linewidth}}}]{
                \texttt{\tiny{[P1$\rangle$GM]}}
                \texttt{Expression: A vertical block of X's in the left three columns, with empty cells in the right two columns.} \\
            }
        }
    }
    & & \\ \\

    \theutterance \stepcounter{utterance}  
    & & & \multicolumn{2}{p{0.3\linewidth}}{
        \cellcolor[rgb]{0.9,0.9,0.9}{
            \makecell[{{p{\linewidth}}}]{
                \texttt{\tiny{[GM$|$GM]}}
                \texttt{Expression: A vertical block of X's in the left three columns, with empty cells in the right two columns.} \\
            }
        }
    }
    & & \\ \\

    \theutterance \stepcounter{utterance}  
    & & \multicolumn{4}{p{0.6\linewidth}}{
        \cellcolor[rgb]{0.9,0.9,0.9}{
            \makecell[{{p{\linewidth}}}]{
                \texttt{\tiny{[GM$\rangle$P2]}}
                \texttt{You are given three grids, where each of them is 5 by 5 in size.} \\
\texttt{Grids have empty cells marked with "▢" and filled cells marked with "X".} \\
\texttt{You are also given a referring expression that describes one of the given grids.} \\
\texttt{Your task is to select the grid that matches the given referring expression.} \\
\texttt{Generate only the number (in text) of the grid that matches the given expression by selecting first, second, or third.} \\
\\ 
\texttt{First:} \\
\\ 
\texttt{□ X X □ X} \\
\texttt{□ X X □ X} \\
\texttt{□ X X □ X} \\
\texttt{□ X X □ X} \\
\texttt{□ X X □ X} \\
\\ 
\texttt{Second:} \\
\\ 
\texttt{□ X X X □} \\
\texttt{□ X X X □} \\
\texttt{□ X X X □} \\
\texttt{□ X X X □} \\
\texttt{□ X X X □} \\
\\ 
\texttt{Third:} \\
\\ 
\texttt{X X X □ □} \\
\texttt{X X X □ □} \\
\texttt{X X X □ □} \\
\texttt{X X X □ □} \\
\texttt{X X X □ □} \\
\\ 
\texttt{Expression: A vertical block of X's in the left three columns, with empty cells in the right two columns.} \\
\texttt{Question: Which grid does the expression refer to?} \\
\texttt{Start with the tag "Answer: ", followed by your selection. Omit any other text.} \\
            }
        }
    }
    & \\ \\

    \theutterance \stepcounter{utterance}  
    & & & \multicolumn{4}{p{0.45\linewidth}}{
        \cellcolor[rgb]{0.95,0.95,0.95}{
            \makecell[{{p{\linewidth}}}]{
                \texttt{\tiny{[P2$|$internal]}}
                \texttt{Reflections from previous episodes:} \\
\texttt{{-}} \\
\texttt{ Failed. The expression was overly complex and repetitive, causing confusion instead of clarity. I tried to parse it literally but got lost in the contradictions and circular descriptions. Next time, I would ignore the convoluted text and simply match the described pattern directly to the visual grids, which would have shown that none of the grids fully matched the described structure.} \\
\texttt{{-}} \\
\texttt{ Answer: First} \\
\\ 
\texttt{Reflection: Succeeded. The expression described first two columns empty, third column filled, last two columns filled. Only Grid 1 matched exactly. I correctly parsed the column structure and compared each grid. Next time, I will continue to systematically verify each column condition against the grids to avoid misreading patterns.} \\
\texttt{{-}} \\
\texttt{ Answer: Second} \\
\\ 
\texttt{Reflection: Succeeded. The expression described a solid vertical block of X's in the right three columns, with left two columns empty. Grid 2 matches exactly: left two columns empty, right three columns filled. I correctly parsed the description and matched it to the visual layout. Next time, I will continue to carefully align column descriptions with the grid orientation.} \\
            }
        }
    }
     \\ \\

    \theutterance \stepcounter{utterance}  
    & & & \multicolumn{4}{p{0.45\linewidth}}{
        \cellcolor[rgb]{0.95,0.95,0.95}{
            \makecell[{{p{\linewidth}}}]{
                \texttt{\tiny{[P2$|$internal]}}
                \texttt{{[}attempt 1 – valid{]}} \\
\texttt{Answer: Third} \\
            }
        }
    }
     \\ \\

    \theutterance \stepcounter{utterance}  
    & & & \multicolumn{4}{p{0.6\linewidth}}{
        \cellcolor[rgb]{0.9,0.9,0.9}{
            \makecell[{{p{\linewidth}}}]{
                \texttt{\tiny{[GM$\langle$P2]}}
                \texttt{Answer: Third} \\
            }
        }
    }
     \\ \\

    \theutterance \stepcounter{utterance}  
    & & & \multicolumn{2}{p{0.3\linewidth}}{
        \cellcolor[rgb]{0.9,0.9,0.9}{
            \makecell[{{p{\linewidth}}}]{
                \texttt{\tiny{[GM$|$GM]}}
                \texttt{third} \\
            }
        }
    }
    & & \\ \\

\end{supertabular}
}

\end{document}
